% Sub-stage F3: the packing of the numerators by stride: the 2-adic classes with their strides and
% residues, the rows of offsets that the packings share, the fields bias, rbainv and dinv of the
% sieve entries, the residues of a denominator for its stride, the bits of the numerators of a
% denominator, the row indices of the first two stages with the offsets and the test small, and the
% stride in the extraction; with the reductions on 64-bit words. Proposed names are set with \lname
% and become \lean once the declarations exist.
\chapter{The packing of the numerators}\label{chap:packing}

F3 models how ratpoints packs the numerators of a denominator $b$ into its bit arrays. The
numerators that the 2-adic information admits for $b$ lie, in most classes of $b \bmod 64$, in a
single residue class $a_0 \bmod 2^k$, and ratpoints sieves only that class: bit $n$ of the sieve of
$b$ stands for the numerator $a_0 + 2^k n$, not for $n$, so that a bit array covers $2^k$ times as
many numerators. The stride $2^k$ and the residue $a_0$ are those of the class of $b \bmod 64$,
which \verb|get_2adic_info| computes once per run with the pattern of the class packed to the
stride. Each modulus then reads, for the denominator $b$, the row of the residue
$b\, 2^{-k} \bmod p$ instead of that of $b \bmod p$, at the bit array
$(i + \mathrm{offset}) \bmod p$ for the word number $i$, with an offset that ratpoints computes once
per run for each packing and each modulus; the first two stages reduce these row indices by a
multiplication with the reciprocal of the modulus where that is exact, and the extraction steps
through the numerators by $2^k$. Until now the model sieved every class with the stride~$1$, the
simplification of Chapter~\ref{chap:denominators}; F3 makes it faithful.

The rows are unchanged when the residue and the numerator are multiplied by the same unit modulo the
modulus, so the packed sieve keeps exactly the numerators of the class $a_0 \bmod 2^k$ that the
unpacked sieve keeps, in the same increasing order, and every numerator outside that class is one
that the 2-adic information rejects anyway. So the output of \verb|findPoints| does not change, nor
do the statements that comparator certifies, \verb|Params| or \verb|Params.WF|: the well-formedness
already makes every modulus odd, above $1$, below $\mathrm{maxPrimeEven}$ and a product of primes of
the table, which is all the packing needs of the parameters. The contracts of the sieve keep their
conclusions with the 2-adic words as ghost arguments, and their hypotheses grow: the class packs its
word, the offsets and the residues are the right ones, and the tables of inverses of the sieve
entries invert. Comparator checks the statements again, since the proofs change.

F3 also makes the reductions by a reciprocal compute on 64-bit words, as the C code does: the
product of the reciprocal and the value is the C's 64-bit product, and the high half of the 128-bit
product, for which Lean core has no type, comes from the 32-bit halves. With the loops over the
moduli stepping two cursors, as the C code steps \verb|ssp| and \verb|mg|, the model with the
packing runs the suite \verb|test4high| about a quarter faster than before F3 (1m32s against 2m02s
on one pinned core).

\section{What the C code does}\label{sec:packingC}

\paragraph{The classes.} For each class $b \bmod 64$ with the word $w$ of the numerators
$a \bmod 64$ that the 2-adic information admits (Definition~\ref{def:twoAdic}),
\verb|get_2adic_info| computes (\csource{find_points.c}{1058}{1104}): with the lowest admitted
numerator $a = \mathrm{ctz}(w)$, the largest $k \le 6$ below $\mathrm{RP\_NUM\_STRIDES} = 7$ such
that every set bit of $w$ is $\equiv a \pmod{2^k}$, by a loop over $k = 1, 2, \dots$ that tests $w$
against the mask of the multiples of $2^k$ shifted by $a \bmod 2^k$ and stops at the first failure,
then decrements; $a_0 = a \bmod 2^k$; and the packed pattern, bit $t$ of which is bit
$(a_0 + 2^k t) \bmod 64$ of $w$, by scattering each set bit $a$ to
$((a - a_0) \mathbin{\&} 63) \gg k$ and doubling the pattern up to 64 bits; for $w = 0$ the class is
empty ($k = 0$, $a_0 = 0$). It sets bit $b$ of the word \verb|den_bits| of the classes that admit a
numerator. The type \verb|rp_num_class| (\csource{rp-private.h}{494}{495}) holds the pattern, $k$,
$a_0$ and a pointer to the class's row of offsets. On every curve $k \le 3$ (an exhaustive search:
all polynomial functions modulo 64 for odd $b$, all truncated reversed coefficients for even $b$);
the code allows $k \le 6$, and the model with it.

\paragraph{The offsets.} \verb|num_packings| (\csource{find_points.c}{1119}{1130}) counts the
distinct keys $2^k + a_0$ of the classes that admit a numerator; the array of offsets has
$\mathrm{num\_packings} \cdot \max(1, \mathrm{sp3\_max})$ entries
(\csource{find_points.c}{5119}{5119}). \verb|class_offsets| (\csource{find_points.c}{1142}{1183},
called at \csource{find_points.c}{5132}{5132}) takes the largest stride $k_{\max}$ of these classes
and for each sieve entry of modulus $p$ two halving chains, $x \mapsto (x + p)/2$ for odd $x$ and
$x/2$ for even $x$, from $1$ and from $\mathrm{rbainv} = 256^{-1} \bmod p$: they give
$\mathrm{dinv}[k] = 2^{-k} \bmod p$, which it stores in the sieve entry, and
$(2^k\, 256)^{-1} \bmod p$ for $k \le k_{\max}$. Then each packing gets one row, shared by the
classes of the same key, with the offset $\mathrm{RP\_MULMOD}(a_0 (2^k\, 256)^{-1}) + \mathrm{bias}$
for each sieve entry, where the bias $p \lceil 2^{31} / p \rceil$
(\csource{find_points.c}{1815}{1815} for a prime, \csource{find_points.c}{2424}{2424} and
\csource{find_points.c}{2456}{2456} for a power and a product) and $\mathrm{rbainv}$
(\csource{find_points.c}{1816}{1816}, \csource{find_points.c}{2425}{2425},
\csource{find_points.c}{2457}{2457}) are fields of the sieve entry, set when it is made, and
$\mathrm{RP\_ROW\_BIAS} = 2^{31}$ (\csource{rp-private.h}{387}{387}) is the lower bound of the bias,
with the static assertion $2\,\mathrm{RP\_ROW\_BIAS} - 1 \le \mathrm{RP\_MULMOD\_LIMIT}$
(\csource{rp-private.h}{426}{427}).

\paragraph{One denominator.} \verb|run_block| takes the class of $b \bmod 64$ and, before each call
of \verb|sift|, \verb|fill_bp_list| (\csource{find_points.c}{3639}{3679}) computes for the moduli of
the first two stages the residues $b\, 2^{-k} \bmod p$ from $\mathrm{dinv}[k]$ by
$\mathrm{RP\_MULMOD}$, in one of three branches chosen by the size of $b$ (the product directly,
after reducing $b$ by $\mathrm{RP\_MULMOD}$, or after reducing it by a division). In \verb|sift|,
\verb|numerators_to_bits| (\csource{find_points.c}{3355}{3370}) turns the interval of numerators
into the bits $\lceil(\mathrm{low} - a_0)/2^k\rceil$ to $\lfloor(\mathrm{high} - a_0)/2^k\rfloor$,
by arithmetic shifts of longs, in a plain form and in a split form that avoids the overflow of the
plain one for heights above $\mathrm{LONG\_MAX} - 64$; for the whole line, the model's setting, the
C code takes the split form exactly for those heights (the test of \verb|exact_bounds| and the path
through \verb|interval_bits_exact|, lines that belong to the search intervals of stage~H). Then,
once an interval of numerators is not empty, \verb|set_rows| (\csource{find_points.c}{3149}{3170})
takes for each modulus the row of that residue and the class's offset, and if the bits are not
empty, they are sifted in pieces with the masks of the ends (\csource{find_points.c}{3563}{3617}).
In \verb|_ratpoints_sift0| the test \verb|small| (\csource{sift.c}{530}{531}) decides whether the
row indices of the piece can be reduced by $\mathrm{RP\_MULMOD}$ with the reciprocals of
\verb|magics|: the first stage computes for each modulus
$r = (-w_{\mathrm{low}} - \mathrm{offset}) \bmod p$ by \verb|mod_mul| (\csource{sift.c}{243}{247})
or by the division \verb|mod| (\csource{sift.c}{370}{389}) and starts at bit array $p - r$ of the
row, which is $\equiv w_{\mathrm{low}} + \mathrm{offset} \pmod p$, and the second stage reads the
bit array $(i + \mathrm{offset}) \bmod p$ in both arms. The extraction gives bit $t$ of word
$\kappa$ of the bit array with the word number $i$ the numerator $a_0 + 2^k (256 i + 64 \kappa + t)$
(\csource{sift.c}{1221}{1222}, \csource{sift.c}{1264}{1264}, \csource{sift.c}{262}{264}). The third
stage does not change: \verb|fill_checks| reduces $b$ itself and tests the numerator itself.

\paragraph{Why the packing changes nothing.} Bit $j < 256$ of the bit array with the word number
$i$, the bit $n = 256 i + j$, stands for $a = a_0 + 2^k n$. Since
$256\,\mathrm{offset}\, 2^k \equiv a_0$ and $\mathrm{bp}\, 2^k \equiv b \pmod p$, the bit that the
row of $\mathrm{bp}$ gives at index $256(i + \mathrm{offset}) + j \equiv a\, 2^{-k}$ is its
admission of $a\, 2^{-k}$, which is the admission of $a$ by the row of $b$: a row is unchanged when
the residue and the numerator are multiplied by the same unit $u$ (for a prime, it reads the table
at $a\, b^{-1}$; for a prime power, at $a\, b^{-1}$ or at $b\, a^{-1}$; both are unchanged), given
that the tables of inverses invert. The packed pattern gives the bit of $w$ at $a \bmod 64$. The
bits from the low to the high end, mapped to $a_0 + 2^k n$, are the numerators of $[-H, H]$ that are
$\equiv a_0 \pmod{2^k}$, in increasing order, and the numerators of $[-H, H]$ outside that class are
rejected by $w$. So the packed sieve gives the numerators that the unpacked sieve keeps, in the same
order: a data refinement.

\section{The constants and the reductions on 64-bit words}\label{sec:words64}

\begin{definition}[The constants of the packing]\label{def:packingConstants}
  \uses{def:rbaPack, def:jacobiTables}
  \lean{Ratpoints.rowBias, Ratpoints.two_mul_rowBias_sub_one_le_mulModLimit, Ratpoints.numStrides, Ratpoints.one_le_numStrides, Ratpoints.rbaShift, Ratpoints.rbaLength_eq_two_pow_rbaShift, Ratpoints.rbaShift_lt_thirty_two}
  \leanok
  The lower bound $\mathrm{rowBias} = 2^{31}$ of the multiple of a modulus in the offset of a row
  (\verb|RP_ROW_BIAS|, \csource{rp-private.h}{387}{387}), with the C's static assertion
  $2\,\mathrm{rowBias} - 1 \le \mathrm{mulModLimit}$ as a fact checked by evaluation
  (\csource{rp-private.h}{426}{427}); the number $\mathrm{numStrides} = 7$ of the strides
  (\verb|RP_NUM_STRIDES|, \csource{rp-private.h}{476}{476}), with the fact
  $1 \le \mathrm{numStrides}$; and $\mathrm{rbaShift} = 8$ (\verb|RBA_SHIFT|,
  \csource{rp-private.h}{126}{126}, \csource{rp-private.h}{173}{173},
  \csource{rp-private.h}{232}{232}, \csource{rp-private.h}{258}{258},
  \csource{rp-private.h}{282}{282}), with the facts $\mathrm{rbaLength} = 2^{\mathrm{rbaShift}}$ and
  $\mathrm{rbaShift} < 32$. No proof uses their values.
\end{definition}

\begin{definition}[The high half of a product]\label{def:mulHi}
  \lean{Ratpoints.Sieve.mulHi}
  \leanok
  The high $64$ bits of the $128$-bit product of two
  $64$-bit words, from the four products of their $32$-bit halves, for the product of type
  \verb|__uint128_t| that gcc provides and Lean core does not (\csource{rp-private.h}{404}{408},
  \csource{rp-private.h}{409}{410}, ranges it shares with Definition~\ref{def:mulMod}). The
  reduction with the reciprocal $m$ (Definition~\ref{def:mulMod}) becomes
  $\mathrm{mulHi}(m u \bmod 2^{64}, p)$, with the C's $64$-bit product of $m$ and $u$, and the
  quotient $\mathrm{mulHi}(u, m)$, the arguments cast to words as the macros cast them to
  \verb|unsigned long|.
\end{definition}

\begin{lemma}[The words compute the products]\label{lem:mulHi}
  \uses{def:mulHi, def:mulMod}
  \lean{Ratpoints.Sieve.toNat_mulHi, Ratpoints.Sieve.mulMod_eq, Ratpoints.Sieve.mulDiv_eq}
  \leanok
  The high half is $\lfloor a b / 2^{64} \rfloor$ for all words
  $a, b$; so the reduction is $\lfloor (m u \bmod 2^{64}) (p \bmod 2^{64}) / 2^{64} \rfloor$ and the
  quotient $\lfloor (u \bmod 2^{64}) (m \bmod 2^{64}) / 2^{64} \rfloor$, and Lemma~\ref{lem:mulMod}
  keeps the statement of the remainder, while the quotient now needs $1 < p$, since the reciprocal
  $2^{64}$ of $1$ wraps to $0$, as in the C code.
\end{lemma}

\begin{proof}
  \leanok
  With $a = 2^{32} a_1 + a_0$ and $b = 2^{32} b_1 + b_0$,
  $a b = 2^{64} a_1 b_1 + 2^{32}(a_1 b_0 + a_0 b_1) + a_0 b_0$; the middle sum of the low halves of
  the cross products and the high half of $a_0 b_0$ is below $3 \cdot 2^{32}$, and its high half is
  the carry.
\end{proof}

\begin{definition}[The remainder of a long of either sign]\label{def:modMul}
  \uses{def:mulMod}
  \lean{Ratpoints.Sieve.modMul}
  \leanok
  For an integer $a$, the reduction of $|a|$ with the
  reciprocal, and $p$ minus it for negative $a$ with a nonzero remainder (\verb|mod_mul|,
  \csource{sift.c}{243}{247}).
\end{definition}

\begin{lemma}[The remainder of a long]\label{lem:modMul}
  \uses{def:modMul, lem:mulMod}
  \lean{Ratpoints.Sieve.modMul_magic}
  \leanok
  For $0 < p$ and $|a|\, p < 2^{64}$,
  $\mathrm{modMul}(a, p, \mathrm{magic}(p)) = a \bmod p$, the Euclidean remainder; \verb|modLong| is
  the Euclidean remainder for a modulus $p \ge 0$ (Definition~\ref{def:useCLong}, whose lemma moves
  to a file without Mathlib).
\end{lemma}

\begin{proof}
  \leanok
  Lemma~\ref{lem:mulMod} for $|a|$; for negative $a$, $p - (|a| \bmod p) = a \bmod p$ when the
  remainder is not $0$.
\end{proof}

\section{The classes}\label{sec:classes}

\begin{definition}[The classes of the numerators]\label{def:numClasses}
  \uses{def:ctz, def:packingConstants}
  \lean{Ratpoints.Sieve.NumClass, Ratpoints.Sieve.classMasks, Ratpoints.Sieve.packPattern, Ratpoints.Sieve.numClasses}
  \leanok
  A class: its packed pattern, a word; its stride exponent $k$ and residue $a_0$; and the position
  of its row in the array of offsets, absent until the offsets are computed
  (\csource{rp-private.h}{494}{495}). The masks of the multiples of $2^k$ for $k \le 6$, the C's
  literal table of seven words, read as $0$ beyond (\csource{find_points.c}{1058}{1062}). The packed
  pattern of a word for $k$ and $a_0$: each set bit $a$ scattered to $((a - a_0) \bmod 64) / 2^k$,
  then the pattern doubled up to $64$ bits (\csource{find_points.c}{1082}{1094}). The classes of the
  $64$ words of the 2-adic information (Definition~\ref{def:twoAdic}) and the word \verb|den_bits|
  (\csource{find_points.c}{995}{995}, \csource{find_points.c}{1064}{1081},
  \csource{find_points.c}{1095}{1097}, \csource{find_points.c}{1099}{1104},
  \csource{find_points.c}{1106}{1114}): for each word, the stride by the loop over the masks with
  its stop and decrement, the residue, and the packed pattern, or the word itself for $k = 0$.
\end{definition}

\begin{lemma}[A class packs its word]\label{lem:numClasses}
  \uses{def:numClasses, def:andRows, lem:twoAdic_point}
  \lean{Ratpoints.Sieve.NumClass.WF, Ratpoints.Sieve.NumClass.Packs, Ratpoints.Sieve.NumClass.numerator, Ratpoints.Sieve.getLsbD_classMasks, Ratpoints.Sieve.packPattern_spec, Ratpoints.Sieve.numClasses_spec, Ratpoints.Sieve.packs_getElem!_fst_run_numClasses, Ratpoints.Sieve.NumClass.Packs.classBit_bits, Ratpoints.Sieve.NumClass.Packs.emod_eq_of_classBit, Ratpoints.Sieve.NumClass.Packs.bits_eq_zero_iff, Ratpoints.Sieve.NumClass.Packs.congr, Ratpoints.Sieve.snd_run_numClasses}
  \leanok
  The class of the word $w$ packs it: $k < \mathrm{numStrides}$, $k \le 6$, $a_0 < 2^k$, every set
  bit $a$ of $w$ is $\equiv a_0 \pmod{2^k}$, and bit $t$ of the pattern is bit
  $(a_0 + 2^k t) \bmod 64$ of $w$; and it has no row of offsets yet. So the pattern's bit for the
  bit number $n$ is the bit of $w$ for the numerator $a_0 + 2^k n$, a numerator of $w$ lies in the
  class $a_0 \bmod 2^k$, and the pattern is $0$ exactly when $w$ is; whether a class packs a word
  depends only on its pattern, its stride and its residue. Bit $b$ of \verb|den_bits| is set exactly
  when the word of $b$ is not $0$, so \verb|den_bits| is the word of Definition~\ref{def:twoAdic}.
\end{lemma}

\begin{proof}
  \leanok
  The mask of $k$ has the bits $\equiv 0 \pmod{2^k}$ for $k \le 6$ and none beyond, which ends the
  stride loop by $k = 7$ at the latest, whatever $\mathrm{numStrides}$ is; an invariant of that loop
  says that the set bits are $\equiv \mathrm{ctz}(w) \pmod{2^{k-1}}$. The scatter puts bit $a$ at
  $(a - a_0)/2^k$, since every set bit is at least $a_0$, and the doubling repeats the pattern with
  the period $64/2^k$, which $2^k \mid 64$ makes the bits $(a_0 + 2^k t) \bmod 64$.
\end{proof}

\section{The offsets of the rows}\label{sec:offsets}

\begin{definition}[The rows of offsets]\label{def:classOffsets}
  \uses{def:numClasses, def:SieveEntry, def:mulMod, def:packingConstants}
  \lean{Ratpoints.Sieve.numPackings, Ratpoints.Sieve.halve, Ratpoints.Sieve.classRows, Ratpoints.Sieve.classOffsets}
  \leanok
  The number of the packings, the distinct keys $2^k + a_0$ of the classes with a nonzero pattern,
  kept in two words as the C code keeps them, and at least $1$
  (\csource{find_points.c}{1119}{1130}); the halving step modulo an odd $p$
  (\csource{find_points.c}{1159}{1160}); the rows: for each class with a pattern, the row of its key
  if it has one, else a new row whose offsets are
  $\mathrm{mulMod}(a_0 \cdot (2^k\,\mathrm{rbaLength})^{-1}) + \mathrm{bias}$ for each sieve entry,
  with the class getting the position of its row (\csource{find_points.c}{1164}{1182}); and
  \verb|class_offsets|: the largest stride of the classes with a pattern, the two halving chains of
  each sieve entry from $1$ and from $\mathrm{rbainv}$ up to it, the first stored as the field
  $\mathrm{dinv}$ of the entry, and the rows (\csource{find_points.c}{1142}{1158},
  \csource{find_points.c}{1161}{1162}, \csource{find_points.c}{1183}{1183}). The bias,
  $\mathrm{rbainv}$ and $\mathrm{dinv}$ are fields of the sieve entries
  (Definition~\ref{def:SieveEntry}).
\end{definition}

\begin{lemma}[The offsets]\label{lem:classOffsets}
  \uses{def:classOffsets, lem:mulMod, lem:numClasses, def:packingConstants, def:barrett}
  \lean{Ratpoints.Sieve.NumClass.key, Ratpoints.Sieve.NumClass.WF.k_eq_and_a0_eq, Ratpoints.Sieve.IsFirstKey, Ratpoints.Sieve.firstKeys, Ratpoints.Sieve.numPackings_spec, Ratpoints.Sieve.classOffset, Ratpoints.Sieve.classRows_spec, Ratpoints.Sieve.halveIter, Ratpoints.Sieve.halveChain, Ratpoints.Sieve.halveIter_mul_two_pow_mod, Ratpoints.Sieve.SieveEntry.Packable, Ratpoints.Sieve.SieveEntry.PackedFor, Ratpoints.Sieve.SieveEntry.PackedFor.congr, Ratpoints.Sieve.OffsetsFor, Ratpoints.Sieve.OffsetsFor.take, Ratpoints.Sieve.mul_mul_lt_two_pow_64, Ratpoints.Sieve.classOffsets_spec}
  \leanok
  The number of the packings is $\max(1, f)$ for the number $f$ of the classes with a pattern whose
  key no earlier class with a pattern has (the classes being well formed, so that the keys are below
  $2^{\mathrm{numStrides}}$). For sieve entries of odd moduli $1 < p \le \mathrm{maxPrimeEven}$ with
  the reciprocal of $p$, $\mathrm{rbainv} < p$ with
  $\mathrm{rbainv} \cdot \mathrm{rbaLength} \equiv 1 \pmod p$ and the bias a multiple of $p$ in
  $[\mathrm{rowBias}, \mathrm{rowBias} + p)$ (\emph{packable}), classes with $k \le 6$,
  $k < \mathrm{numStrides}$ and $a_0 < 2^k$, and an array of at least $f$ times the number of
  entries: the classes keep their patterns, strides and residues, the entries change only in
  $\mathrm{dinv}$, and each class with a pattern has a row at a position $\mathrm{pos}$ whose offset
  $o$ for the entry $n$ satisfies $\mathrm{rowBias} \le o < \mathrm{rowBias} + 2p$ and
  $o\, 2^k\,\mathrm{rbaLength} \equiv a_0 \pmod p$, with
  $\mathrm{dinv}[k] \cdot 2^k \equiv 1 \pmod p$ and $\mathrm{dinv}[k] < p$.
\end{lemma}

\begin{proof}
  \leanok
  The halving step multiplies by $2^{-1}$ modulo an odd $p$ and stays below $p$. The rows occupy the
  disjoint ranges $[r\, n_p, (r + 1)\, n_p)$, for the number $n_p$ of the sieve entries, one for
  each class with a pattern whose key no earlier class with a pattern has, and the keys $2^k + a_0$
  with $a_0 < 2^k$ determine $(k, a_0)$; an offset depends only on the key. The halvings of
  $\mathrm{rbainv} < p$ stay below $p$, so $a_0 (2^k\,\mathrm{rbaLength})^{-1} < 64\, p$, which
  times $p$ is below $2^{64}$: its reduction is exact by Lemma~\ref{lem:mulMod}.
\end{proof}

\section{The residues of a denominator}\label{sec:bp}

\begin{definition}[The residues of a denominator for its stride]\label{def:fillBpList}
  \uses{def:SieveEntry, def:mulMod, def:useCLong, def:jacobiTables, def:barrett}
  \lean{Ratpoints.Sieve.fillBpList}
  \leanok
  For the denominator $b$, the stride exponent $k$, the sieve entries of the first two stages and
  their reciprocals: the residues $b\,\mathrm{dinv}[k] \bmod p$, by $\mathrm{mulMod}$ of the product
  for $b \le \mathrm{mulModLimit}/\mathrm{maxPrimeEven}$, of the reduced $b$ times
  $\mathrm{dinv}[k]$ for $b \le \mathrm{mulModLimit}$, and of $(b \bmod p) \,\mathrm{dinv}[k]$ by
  \verb|modLong| above, the branch chosen once for $b$ (\csource{find_points.c}{3639}{3679}).
\end{definition}

\begin{lemma}[The residues]\label{lem:fillBpList}
  \uses{def:fillBpList, lem:mulMod}
  \lean{Ratpoints.Sieve.fillBpList_eq_map, Ratpoints.Sieve.mulMod_mul_magic, Ratpoints.Sieve.mulMod_mulMod_mul_magic, Ratpoints.Sieve.mulMod_toNat_modLong_mul_magic, Ratpoints.Sieve.mod_le_two_pow_10}
  \leanok
  For entries with $0 < p \le \mathrm{maxPrimeEven}$ and $\mathrm{dinv}[k] < p$ and their
  reciprocals, the residues are $b\, \mathrm{dinv}[k] \bmod p$, in all three branches.
\end{lemma}

\begin{proof}
  \uses{lem:classOffsets, lem:packingLongs}
  \leanok
  Each value $u$ that is reduced has $u\, p < 2^{64}$:
  $b\,\mathrm{dinv}[k] \le \mathrm{mulModLimit}$ in the first branch
  (Lemma~\ref{lem:packingLongs}), $b \le \mathrm{mulModLimit}$ for the first reduction of the
  second, and $(b \bmod p)\,\mathrm{dinv}[k] < p^2$ for the last reduction of the second and of the
  third, which times $p$ is below $2^{64}$ as the products of the offsets are
  (Lemma~\ref{lem:classOffsets}), since $\mathrm{mulModLimit} < 2^{32}$ and $p \le 2^{10}$;
  Lemma~\ref{lem:mulMod}, and \verb|modLong| is the Euclidean remainder
  (Definition~\ref{def:useCLong}).
\end{proof}

\section{The rows under a unit}\label{sec:scaling}

\begin{lemma}[A row is unchanged under a unit]\label{lem:admits_mul}
  \uses{def:SieveEntry, def:examinePower, def:admits}
  \lean{Ratpoints.Sieve.Factor.HasInverses, Ratpoints.Sieve.Entry.admits_mul, Ratpoints.Sieve.Power.admits_mul, Ratpoints.Sieve.Factor.HasInverses.admits_mod, Ratpoints.Sieve.Factor.HasInverses.admits_mul, Ratpoints.Sieve.SieveEntry.admits_mul, Ratpoints.Sieve.Entry.admits_mod_of_dvd, Ratpoints.Sieve.Power.admits_mod_of_dvd}
  \leanok
  If the modulus of each factor of a sieve entry divides $p$ and its table of inverses inverts its
  units (the nonzero residues of a prime, the residues prime to $\ell$ of a prime power $\ell^e$),
  and $v u \equiv 1 \pmod p$, then the row of $u b$ admits $u a$ exactly when the row of $b$ admits
  $a$; a row of a factor depends only on the residues modulo a multiple of its modulus.
\end{lemma}

\begin{proof}
  \leanok
  The table of a prime is read at $a\, \mathrm{inv}(b)$, that of a prime power at $a\,
  \mathrm{inv}(b)$ or at $b\,\mathrm{inv}(a)$; with $u b$ and $u a$ the index is $u a\,
  \mathrm{inv}(u b) \equiv u a\, \mathrm{inv}(u b)\, b\, \mathrm{inv}(b) \equiv a\,
  \mathrm{inv}(b)$, using only that the tables invert and that $u$ has the inverse $v$, so no
  primality and no field of residues; the tests for $0$ of a prime and for divisibility by the prime
  $\ell$ of a prime power $\ell^e$ are unchanged, since $u$ is a unit modulo the modulus of the
  factor, which divides $p$, and so modulo $\ell$.
\end{proof}

\begin{lemma}[The rows read in the packing]\label{lem:reads_row}
  \uses{lem:admits_mul, def:rows, def:andRows}
  \lean{Ratpoints.Sieve.SieveSpec.Reads, Ratpoints.Sieve.SieveSpec.getLsbD_wordAt_mk, Ratpoints.Sieve.SieveEntry.reads_row, Ratpoints.Sieve.getLsbD_andWords}
  \leanok
  A sieve spec reads a test $f$ of the bit numbers when bit $t$ of word $y$ of the piece from the word
  number $i$ is $f(\mathrm{rbaLength}\, i + 64 y + t)$ for $t < 64$. For a sieve entry with $0 < p$
  and tables that invert, $v\, 2^k \equiv 1$, $o\, 2^k\,\mathrm{rbaLength} \equiv a_0$ and
  $2^k\, \mathrm{bp} \equiv b \pmod p$, the spec of the row of $\mathrm{bp}$ at the offset $o$ reads
  the admission of $a_0 + 2^k j$ by the row of $b$; the words of a piece and-ed over specs that read
  tests read the conjunction of the tests and the class word.
\end{lemma}

\begin{proof}
  \leanok
  The row bit at the index $\mathrm{rbaLength}(i + o) + 64 y + t$ admits that index for
  $\mathrm{bp}$ (Definitions~\ref{def:row} and~\ref{def:rows}); multiplied by the unit $2^k$ with
  the inverse $v$ (Lemma~\ref{lem:admits_mul}) it is the admission of
  $2^k(\mathrm{rbaLength}\, i + 64 y + t) + a_0$ for $b$.
\end{proof}

\section{The row indices of the first two stages}\label{sec:rowIndices}

\begin{lemma}[The row indices]\label{lem:rowIndices}
  \uses{lem:mulMod, lem:modMul, def:packingConstants, def:rows, def:barrett, def:useCLong}
  \lean{Ratpoints.Sieve.SieveSpec.Fits, Ratpoints.Sieve.SmallFits, Ratpoints.Sieve.smallFits_decide, Ratpoints.Sieve.SmallFits.mono, Ratpoints.Sieve.SieveSpec.Fits.natAbs_neg_sub_le, Ratpoints.Sieve.SieveSpec.Fits.add_nonneg_and_toNat_le, Ratpoints.Sieve.SieveSpec.Fits.ite_modMul_modLong_eq_emod, Ratpoints.Sieve.SieveSpec.Fits.ite_mulMod_modLong_eq_toNat_emod, Ratpoints.maxPrimeEven_pos}
  \leanok
  A sieve spec \emph{fits} when $0 < p \le \mathrm{maxPrimeEven}$ and
  $\mathrm{rowBias} \le \mathrm{offset} < \mathrm{rowBias} + 2p$; the test \verb|small| of a piece
  from $w_0$ to $w_0 + n$ gives $-\mathrm{rowBias} \le w_0$ and
  $w_0 + n \le \mathrm{rowBias} - 2\,\mathrm{maxPrimeEven}$. Then the remainder $r$ with which the
  first stage starts each modulus, at the bit array $p - r$ of its row, computed by \verb|mod_mul|
  under \verb|small| and by \verb|mod| otherwise, is $(-w_0 - \mathrm{offset}) \bmod p$, and the
  second stage's index for a word number $i$ of the piece, computed by $\mathrm{mulMod}$ under
  \verb|small| and by \verb|mod| otherwise, is $(i + \mathrm{offset}) \bmod p$.
\end{lemma}

\begin{proof}
  \leanok
  Under \verb|small| the values have absolute value at most
  $2\,\mathrm{rowBias} - 1 \le \mathrm{mulModLimit} < 2^{32}$, and $p \le 2^{10}$, so the reductions
  are exact (Lemmas~\ref{lem:mulMod} and~\ref{lem:modMul}); \verb|mod| is the Euclidean remainder.
  The proof uses the static assertion, never the value of $\mathrm{rowBias}$.
\end{proof}

The model passes the reciprocals of the first two stages to the stages as \verb|magics| holds them,
from the position $0$ for the first stage and from $s_1$ for the second, and the loops over the
moduli step the sieve specs and the reciprocals together, by two cursors, as the C code steps
\verb|ssp| and \verb|mg|. The first stage (Definition~\ref{def:andRows}), the second
(Definition~\ref{def:secondStage}) and one piece (Definition~\ref{def:siftPiece}), which computes
\verb|small|, take the offsets of the sieve specs (Definition~\ref{def:rows}); their contracts take
fitting specs, the reciprocals of their moduli and the bounds of \verb|small|, with their
conclusions in the words of Lemma~\ref{lem:reads_row}.

\section{The bits of one denominator}\label{sec:bits}

\begin{definition}[The bits of an interval of numerators]\label{def:numeratorsToBits}
  \uses{def:numClasses}
  \lean{Ratpoints.Sieve.numeratorsToBits}
  \leanok
  For the ends of an interval of numerators, $64$-bit signed words, and a class: the bits
  $-((a_0 - \mathrm{low}) \gg k)$ and $(\mathrm{high} - a_0) \gg k$ in the plain form, and
  $(\mathrm{low} \gg k) + (((\mathrm{low} \mathbin{\&} (2^k - 1)) - a_0 + 2^k - 1) \gg k)$ and
  $(\mathrm{high} \gg k) + (((\mathrm{high} \mathbin{\&} (2^k - 1)) - a_0) \gg k)$ in the split
  form, with the arithmetic shift and the and of two's complement, as the C code computes them on
  longs (\csource{find_points.c}{3355}{3370}).
\end{definition}

\begin{lemma}[The bits of the heights]\label{lem:numeratorsToBits}
  \uses{def:numeratorsToBits, lem:numClasses}
  \lean{Ratpoints.Sieve.toInt_numeratorsToBits, Ratpoints.Sieve.toInt_numeratorsToBits_height, Ratpoints.Sieve.le_numerator_and_numerator_le_iff, Ratpoints.Sieve.map_range_numerator_eq_filter_numerators, Ratpoints.Sieve.toInt_numeratorsToBits_false, Ratpoints.Sieve.toInt_numeratorsToBits_true, Ratpoints.Sieve.filter_numerators_emod_eq}
  \leanok
  For a class with $k \le 6$, $k < \mathrm{numStrides}$ and $a_0 < 2^k$, the split form gives
  $-\lfloor(a_0 - \mathrm{low})/2^k\rfloor$ and $\lfloor(\mathrm{high} - a_0)/2^k\rfloor$ without
  overflow, and so does the plain form when $a_0 - \mathrm{low} < 2^{63}$ and
  $\mathrm{high} - a_0 \ge -2^{63}$; for the ends $\mp H$ with $H < 2^{63}$ and the form the C code
  takes, the bits are $-\lfloor (H + a_0)/2^k\rfloor$ and $\lfloor (H - a_0)/2^k\rfloor$. A
  numerator $a_0 + 2^k t$ lies in $[\mathrm{low}, \mathrm{high}]$ exactly when $t$ lies between the
  bits, and the bits from the low to the high end, mapped to $a_0 + 2^k t$, are the numerators of
  $[-H, H]$ that are $\equiv a_0 \pmod{2^k}$, in increasing order.
\end{lemma}

\begin{proof}
  \leanok
  The arithmetic shift of a word is the floor division of its integer by $2^k$, the and with
  $2^k - 1$ its Euclidean remainder; each sum stays in the range of a word. The enumeration counts
  the members of a residue class in an interval.
\end{proof}

\begin{definition}[The bits of one denominator]\label{def:siftBits}
  \uses{def:siftPiece, def:rows, def:numClasses, def:rbaPack}
  \lean{Ratpoints.Sieve.setRows, Ratpoints.Sieve.siftBits}
  \leanok
  The sieve specs of a denominator: for each modulus of the first two stages, the row of its residue for the stride
  and the offset of the class's row (\verb|set_rows|, \csource{find_points.c}{3149}{3170}); and the
  bits from $\mathrm{low}$ to $\mathrm{high}$: nothing if $\mathrm{low} > \mathrm{high}$, else the
  pieces of the word numbers $\lfloor \mathrm{low}/\mathrm{rbaLength}\rfloor$ to
  $\lfloor \mathrm{high}/\mathrm{rbaLength}\rfloor$ with the masks
  $\mathrm{low} - \mathrm{rbaLength}\, w_{\mathrm{low}}$ and
  $\mathrm{rbaLength} - 1 - (\mathrm{high} \bmod \mathrm{rbaLength})$
  (\csource{find_points.c}{3563}{3612}). One denominator (Definition~\ref{def:sift}) takes the
  class, the residues and the array of offsets, makes the specs, turns $[-H, H]$ into bits and sifts
  them.
\end{definition}

\begin{lemma}[The bits of one denominator]\label{lem:siftBits}
  \uses{def:siftBits, lem:reads_row, lem:numeratorsToBits, lem:classOffsets, def:siftPiece,
    lem:numClasses, lem:rowIndices}
  \lean{Ratpoints.Sieve.siftBits_spec, Ratpoints.Sieve.packedReads, Ratpoints.Sieve.setRows_map, Ratpoints.Sieve.flatMap_filter_bits_eq, Ratpoints.Sieve.reads_of_mem_packedReads, Ratpoints.Sieve.all_packedReads, Ratpoints.Sieve.fits_of_mem_packedReads, Ratpoints.Sieve.map_magic_packedReads, Ratpoints.Sieve.flatMap_acceptedPoints_eq_pointsOfDen, Ratpoints.Sieve.SieveEntry.PackedFor.fits, Ratpoints.Sieve.NumClass.Packs.filter_filter_emod}
  \leanok
  The bits of an interval give the points of the bits from $\mathrm{low}$ to $\mathrm{high}$ that pass the
  class pattern and the tests that the specs read, none for an empty interval; the specs that
  \verb|set_rows| makes for entries with their residues $b\,\mathrm{dinv}[k] \bmod p$ and the
  offsets of the class read the admissions of the numerators $a_0 + 2^k j$ by the rows of $b$; and
  the bits of $[-H, H]$ in the packing of a class of the word $w$ that pass its pattern give the
  numerators of $[-H, H]$ that pass $w$.
\end{lemma}

\begin{proof}
  \uses{def:checkNum, lem:accepted, lem:staged}
  \leanok
  An invariant of the loop over the pieces, as before over the bits instead of the numerators
  (Definition~\ref{def:sift}); Lemma~\ref{lem:reads_row} with the offsets of
  Lemma~\ref{lem:classOffsets} and the residues of Lemma~\ref{lem:fillBpList};
  Lemmas~\ref{lem:numeratorsToBits} and~\ref{lem:numClasses} for the class.
\end{proof}

\section{The sieve in the packing}\label{sec:packedSieve}

The contracts below one denominator stay on the bit numbers: the walk through the set bits of a word
(Definition~\ref{def:walkWord}) takes the step $2^k$; a bit array through the second stage and the
extraction (Definition~\ref{def:siftBitArray}), the scan and one piece take the class and give the
points of the numerators $a_0 + 2^k j$ of the bits $j$ that survive, quantified over the tests that
the specs read. The contract of one denominator (Definition~\ref{def:sift}) and the lemma that the
words compute the Booleans (Lemma~\ref{lem:sift_eq_checkDen}) keep their conclusions, the points of
the numerators of $[-H, H]$ that pass the 2-adic word $w$ and survive all the entries, with $w$ as a
ghost argument, under the hypotheses that the class packs $w$, has a row at the position
$\mathrm{pos}$, the specs are those of the rows of the entries, the reciprocals those of their
moduli, the residues $b\,\mathrm{dinv}[k] \bmod p$, the tables of inverses invert, and each entry is
packed for the class at its offset. The run (Definition~\ref{def:sieveRun}) takes the classes and
the array of offsets, calls \verb|fill_bp_list| before each denominator, and keeps its conclusion
with the 2-adic words as a ghost argument packed by the classes, under the tables of inverses of the
first $s_2$ entries, their offsets, and \verb|den_bits| matching the classes (the shape of the
forbidden divisors tested with words reads \verb|den_bits|, the others the class's pattern,
Definition~\ref{def:Shape}). The search (Definition~\ref{def:searchPoints}) makes the classes after
the 2-adic information, sizes the array of offsets and computes them after \verb|sieving_info|,
which gives packable entries with their tables of inverses (Lemma~\ref{lem:sievingInfo}).

\begin{lemma}[What the run needs of the search]\label{lem:soundRun}
  \uses{lem:sievingInfo, lem:classOffsets, lem:numClasses, lem:admits_mul, def:stage3, def:twoAdic}
  \lean{Ratpoints.Sieve.IsSoundRun, Ratpoints.Sieve.IsSoundRun.bits_bne_zero, Ratpoints.Sieve.IsSoundRun.getLsbD_db, Ratpoints.Sieve.NumClass.Packs.bits_bne_zero}
  \leanok
  The sieve entries, reciprocals, classes, offsets
  and \verb|den_bits| that the search hands to the run are sound: sound rows, prime entries in the
  third stage, the reciprocals in their order, tables of inverses and the offsets of the classes for
  the first two stages, classes that pack the 2-adic words, and \verb|den_bits| the word of
  Definition~\ref{def:twoAdic}; so a pattern is nonzero exactly when the word of its class is, and
  \verb|den_bits| has the bits of the classes with a pattern. The search establishes it for its
  arguments in its proof.
\end{lemma}

\begin{proof}
  \leanok
  Lemma~\ref{lem:sievingInfo}, Lemma~\ref{lem:numClasses} and Lemma~\ref{lem:classOffsets}; the
  entries that \verb|class_offsets| returns differ only in $\mathrm{dinv}$, which the soundness, the
  shape of a prime entry and the tables of inverses do not read.
\end{proof}

\begin{lemma}[The longs of the packing]\label{lem:packingLongs}
  \uses{lem:classOffsets, def:packingConstants, def:useCLong}
  \lean{Ratpoints.Sieve.NumClass.fitsLong_numerator, Ratpoints.Sieve.neg_two_pow_le_rbaLength_mul_ediv, Ratpoints.Sieve.rbaLength_mul_ediv_add_one_le_two_pow, Ratpoints.Sieve.fitsLong_neg_sub_and_fitsLong_add, Ratpoints.Sieve.SieveEntry.PackedFor.lt_two_pow_32, Ratpoints.Sieve.mul_le_mulModLimit_of_le_div, Ratpoints.Sieve.mul_lt_two_pow_63, Ratpoints.rowBias_le_two_pow_31, Ratpoints.Sieve.natCast_rbaLength_dvd_two_pow, Ratpoints.Sieve.sixty_four_le_rbaLength}
  \leanok
  The values that the C code holds in longs fit a long ($\mathrm{fitsLong}$,
  Definition~\ref{def:useCLong}) for $H < 2^{63}$: the numerator of every bit of the bit arrays of
  $[-H, H]$ (Lemma~\ref{lem:numeratorsToBits}) and the base $a_0 + i\,\mathrm{rbaLength}\, 2^k$ of
  the extraction with its steps, $\mathrm{rbaLength}\, w_{\mathrm{low}}$ and $w_{\mathrm{high}}$,
  $-w_{\mathrm{low}} - \mathrm{offset}$ and $i + \mathrm{offset}$ outside \verb|small|, the offsets,
  below $2^{32}$, and the products of \verb|fill_bp_list|, at most $\mathrm{mulModLimit}$ in the
  first branch, and of \verb|class_offsets|; for $k \le 6$, the multiple of $\mathrm{rbaLength}$ at
  or below a bit $\ge -2^{63-k}$ and the one above a bit $< 2^{63-k}$ lie in
  $[-2^{63-k}, 2^{63-k}]$.
\end{lemma}

\begin{proof}
  \leanok
  $\mathrm{rbaLength}\, 2^k$ is a power of $2$ that divides $2^{63}$
  ($\mathrm{rbaLength} = 2^{\mathrm{rbaShift}}$, $\mathrm{rbaShift} + k \le 63$). So the multiples
  of it at the two ends of the bit arrays of $[-H, H]$, at or below the low end and above the high
  end, lie in $[-2^{63}, 2^{63}]$, and the numerators $a_0 + 2^k j$ of the bits between them, with
  $a_0 < 2^k$, lie in $[-2^{63}, 2^{63})$, which also keeps in range the base of the extraction,
  stepped before each word after the first. The others are bounded by the word numbers, the moduli
  and the bias.
\end{proof}

\section{The command-line program and the differential tests}\label{sec:packingCli}

The program does not change. The differential tests compare every point list in its order and cannot
see the packing, since the output is the same at every stride; the evidence that the packing is that
of ratpoints is the comparison of the model's classes, offsets, residues, bit ranges, masks and test
\verb|small| with those that an instrumented build of the C program prints
(Section~\ref{sec:packing-evidence}), and \verb|#guard| examples against values of the C program:
the classes of words against the C's block for every stride, the packing of the curve
$y^2 = x^3 - 2$ that the C program prints in its \verb|DEBUG| build (classes, \verb|den_bits|, seven
packings, biases, residues and offsets), the rows of offsets against their rule for moduli of every
kind up to $255$, the residues in the three branches, the bits in both forms near $2^{63}$ (the
plain form wrapping at $H = 2^{63} - 1$, where the C code takes the split form), \verb|mod_mul|
against the Euclidean remainder, the high half against the product of natural numbers, and one
denominator in the packing against the reference for both values of the flag of the reversal and at
the heights $1$ to $3$, where a bit interval can be empty or have a negative end.

\section{Deviations}\label{sec:packingDeviations}

Each component differs from the C code in a stated way, recorded in its \verb|Deviation:| paragraph:
\begin{itemize}
  \item the high half of the $128$-bit product from the $32$-bit halves, for the product of type
    \verb|__uint128_t|; natural numbers for the \verb|long| arguments and results of the reductions,
    cast to words as the macros cast them;
  \item the packed pattern of a class a component of its own, called by the loop over the classes,
    as the loop of \verb|relprime| is; the pattern one word, which the first stage writes into every
    word of a bit array (\verb|RBA(packed)|); the class's row of offsets by its position in the
    array, for the pointer; the masks of the strides read as $0$ beyond the table, as the C code
    reads the zero-initialized entries of a larger \verb|RP_NUM_STRIDES|;
  \item the halving step, the loop over the classes of \verb|class_offsets| and the bits of one
    interval of \verb|sift| as functions of their own; one definition for the two copies of
    \verb|mod|, without the hooks of the counts of \verb|sift.c|; the two ends of
    \verb|numerators_to_bits| returned as a pair, for its pointer arguments;
  \item the rows of offsets with an offset for each sieve entry of the run, those of its three
    stages, where ratpoints sizes the rows for the $\mathrm{sp3\_max}$ entries with the reserve that
    the correction of the moduli may take (stage~I); the entries returned by \verb|class_offsets|
    with their field $\mathrm{dinv}$, for the writes into the sieve entries; $\mathrm{dinv}$ an
    array of the strides up to the largest, where the C code has \verb|RP_NUM_STRIDES| entries and
    leaves those above the largest unset;
  \item the fields $\mathrm{bias}$ and $\mathrm{rbainv}$ of a prime's sieve entry set where the
    model makes the entry of a modulus, as its reciprocal is, where \verb|examine_prime| sets them;
  \item the array of the offsets and the residues of the denominator as arguments of one
    denominator, for the pointers of the class and of the arguments; the residues as a list made
    before the call, as \verb|run_block| fills \verb|bp_list|; no measurement of the time of
    \verb|fill_bp_list| (\verb|RP_BP_TIC|, \verb|RP_BP_TOC|);
  \item the bits of the interval on $64$-bit signed words in both forms, with the form that the C
    code takes for the whole line, decided by lines of the search intervals (stage~H); the bit range
    after it, the masks and the extraction base as integers, with the floor division and the
    Euclidean remainder for the shift and the and of a negative long, and their range obligations as
    lemmas (Lemma~\ref{lem:packingLongs});
  \item the reciprocals of the first two stages as lists stepped with the sieve specs, for the
    pointer \verb|mg| into \verb|magics|; \verb|small| a Boolean, for the \verb|int|;
  \item the sieve specs, and with them the rows of the residues, made before the bits of the
    interval, where \verb|sift| turns the interval into bits first and sets up once an interval of
    numerators is not empty (the flag \verb|do_setup|), which for the whole line is at once and so
    before the test $\mathrm{low} \le \mathrm{high}$ of the bits.
\end{itemize}
The lines of the packing that are neither modelled nor left to a later stage are excluded in
\verb|correspondence/ratpoints.toml|: none beyond those already excluded (the timing lines of
\verb|fill_bp_list| are owned with a deviation, as those of \verb|sift| are). Left to stage~F4 are
the start pointers of the chunked first stage (\csource{sift.c}{596}{606}), which read \verb|small|
and \verb|mod_mul| too, and the ends of the rows; to stage~H the search intervals, among them the
choice of the form of \verb|numerators_to_bits| and the clipping of the interval; to stage~I the
reserve of the offsets for the correction of the moduli and the cost model of the packing
(\verb|clipped_padding|).

\section{What changes in the earlier chapters}\label{sec:packingEarlier}

\begin{itemize}
  \item Chapter~\ref{chap:sieve}: the run (Definition~\ref{def:sieveRun}) as in
    Section~\ref{sec:packedSieve}, its node using Definition~\ref{def:fillBpList}, and the reduction
    of the run to the reference (Theorem~\ref{thm:run_reference}) with the classes, the residues and
    the offsets, under the tables of inverses and the offsets of the first $s_2$ entries and
    \verb|den_bits| matching the classes in place of the positivity of the moduli, and with them the
    passages on the run after Lemma~\ref{lem:sift_eq_checkDen}, after Theorem~\ref{thm:refusals} and
    in Section~\ref{sec:staged}; the record of comparator after Theorem~\ref{thm:findPoints}; the
    rules of the translation name \verb|numeratorsToBits| on $64$-bit signed words and the high half
    on $64$-bit words next to the test for common factors; the list of deviations loses its items on
    the residues of $b / 2^k$, on every numerator for the packing and on the plain layout, keeping
    the search region (H) and the correction between the blocks (I).
  \item Chapter~\ref{chap:words}: the sieve specs gain the offset (Definition~\ref{def:rows}); the
    first stage (Definition~\ref{def:andRows}) takes the reciprocals and \verb|small| and computes
    for each modulus $r = (-w_0 - \mathrm{offset}) \bmod p$, starting at bit array $p - r$ of its
    row; the walk (Definition~\ref{def:walkWord}) takes the step; the scan and one piece
    (Definitions~\ref{def:scanPiece} and~\ref{def:siftPiece}) take the class and the reciprocals of
    their moduli, the scan also \verb|small|, which one piece computes; one denominator
    (Definition~\ref{def:sift}) takes the class, the reciprocals, the residues and the offsets, with
    \verb|set_rows| and the bits as components of their own, its node using
    Definitions~\ref{def:siftBits} and~\ref{def:numeratorsToBits} and giving
    \csource{find_points.c}{3149}{3170} and \csource{find_points.c}{3563}{3612} to
    Definition~\ref{def:siftBits}, keeping \csource{find_points.c}{3613}{3617}; their contracts and
    Lemma~\ref{lem:sift_eq_checkDen} as in Section~\ref{sec:packedSieve};
    Definition~\ref{def:rbaPack} gives the lines of \verb|RBA_SHIFT| to
    Definition~\ref{def:packingConstants}, and its proofs use
    $\mathrm{rbaLength} = 2^{\mathrm{rbaShift}}$ and $\mathrm{rbaShift} < 32$ besides
    $\mathrm{rbaPack} > 0$, which hold for every build; the paragraph after
    Lemma~\ref{lem:sift_eq_checkDen} no longer has \verb|sift| compute $b \bmod p$, the run
    computing the residues of $b / 2^k$ before each call, as \verb|fill_bp_list| does; the
    introduction and the list of deviations no longer name the packing as later work (the plain
    layout, the row at index $w \bmod p$); the contract of the trailing zeros and the class word
    move to a new file below \verb|Words.lean|.
  \item Chapter~\ref{chap:setup}: \verb|modLong| (Definition~\ref{def:useCLong}) also models the
    \verb|mod| of \verb|sift.c| (\csource{sift.c}{370}{389}, the same function as in
    \verb|find_points.c| but for the hooks of its counts), and its lemma moves to
    \verb|Ratpoints/Proof/Reductions.lean|; the search (Definition~\ref{def:searchPoints}) takes the
    classes and the offsets (\csource{find_points.c}{5119}{5119},
    \csource{find_points.c}{5132}{5132}); the deviation on \verb|examine_prime| says that its bias
    and \verb|rbainv| are set in the sieve entry that \verb|makeModulus?| makes, as its reciprocal
    is; the record of \verb|find_points_work_1| lacks only what belongs to stages~F4, H and~I.
  \item Chapter~\ref{chap:denominators}: the 2-adic information (Definition~\ref{def:twoAdic}) gives
    the lines of \verb|den_bits| and of the output under \verb|DEBUG| to the classes, and no longer
    keeps the stride~$1$ or computes \verb|den_bits| itself; the shapes (Definition~\ref{def:Shape})
    test the class's pattern, with \lname{Ratpoints.Sieve.Shape.bits\_ne\_zero\_of\_mem\_dens} in
    the node of Lemma~\ref{lem:shape_dens}; the introduction loses the simplification of the
    stride~$1$; the paragraph after Lemma~\ref{lem:twoAdic_point} describes the classes of
    Definition~\ref{def:numClasses} in place of the class word and loses its links to the lines of
    the packing; the deviation on the stride~$1$ and the packing among the work left to later stages
    go.
  \item Chapter~\ref{chap:tables}: the sieve entries (Definition~\ref{def:SieveEntry}) gain
    $\mathrm{bias}$, $\mathrm{rbainv}$ and $\mathrm{dinv}$; the sieve entry of a modulus
    (Definition~\ref{def:makeModulus}) sets the first two with the C's formulas and leaves
    $\mathrm{dinv}$ empty for the rows of offsets (Definition~\ref{def:classOffsets}) to fill; its
    contract (Lemma~\ref{lem:sieveEntry_admits}, with \lname{Ratpoints.IsModulusEntry} and
    \lname{Ratpoints.Sieve.Factor.IsSound}) describes the entry of a modulus that satisfies the
    conditions of the well-formedness (sound rows, tables of inverses, packable), and that of
    \verb|sieving_info| (Lemma~\ref{lem:sievingInfo}) passes this on; the lemmas of a factor's rows
    modulo a multiple move to Lemma~\ref{lem:admits_mul}; the bounds
    $\mathrm{maxPrimeEven} \le 2^{10}$ and $\le 2^{32}$ move to
    \verb|Ratpoints/Proof/Reductions.lean|; the deviations on the fields of the sieve entry that
    serve the packing and the sentence on the first two stages reading the reciprocals with the
    packing say what F3 does.
  \item Chapter~\ref{chap:stages}: the reductions (Definition~\ref{def:mulMod}) through the high
    half, the node using Definition~\ref{def:mulHi}, and Lemma~\ref{lem:mulMod}, whose proofs rest
    on Lemma~\ref{lem:mulHi}, which its node uses, with the quotient for $1 < p$; the shape of a
    prime entry stated field by field, whatever its fields of the packing
    (Definition~\ref{def:stage3}); Lemma~\ref{lem:stage3} states its lemmas over the single factor
    of an entry, names the existing \verb|Ratpoints.Sieve.SieveEntry.IsPrimeEntry.admits_checkOf| in
    place of \verb|Ratpoints.Sieve.CheckSpec.admits_checkOf|, which is folded into it, and
    \lname{Ratpoints.Sieve.bias\_checkOf\_ne\_zero}, and gives \verb|isPrimeEntry_examinePrime| for
    any entry with the fields of a prime; the second stage (Definition~\ref{def:secondStage},
    Lemma~\ref{lem:secondStage}, with \lname{Ratpoints.Sieve.SecondStageWordInvariant}) with the
    reciprocals, \verb|small| and the offsets; a bit array through the second stage and the
    extraction (Definition~\ref{def:siftBitArray}) with the class, the walk of the whole bit array
    stepping its base before each word after the first, as the C code does; the deviations on the
    exact $128$-bit products, on the reciprocal that does not wrap at $p = 1$ and on the row at
    index $i \bmod p$, and the work left to stage~F3, are rewritten or go.
  \item The stage list of the introduction names the reductions by a reciprocal on $64$-bit words in
    the item of F3 and this chapter (Chapter~\ref{chap:packing}), which is input after
    Chapter~\ref{chap:stages}; Chapter~\ref{chap:sieve} already counts the packing among the data
    refinements. The files sections, the uses and the generated table follow.
\end{itemize}

\section{Files, names and order of work}\label{sec:packingFiles}

The constants go into \verb|Ratpoints/Model/Constants.lean|; the high half and \verb|mod_mul| into
\verb|Ratpoints/Model/Tables.lean|; the classes, the rows of offsets and the residues into a new
file \verb|Ratpoints/Model/Packing.lean|, which the run and the search import; the class, the bits,
the specs of a denominator and the changed components of the sieve stay in
\verb|Ratpoints/Model/Sieve.lean|, which gives the points of a reversed curve, the trailing zeros of
a word, the test of a candidate and the exact check (the test for common factors, the third stage,
the check of a numerator) and the points at infinity, about $420$ lines with their guards, to a new
file \verb|Ratpoints/Model/Candidate.lean| below it, so that both stay below $1500$ lines
(\verb|Ratpoints/Model/Sieve.lean| has $1506$ in the prototype). The proofs: a new file without
Mathlib \verb|Ratpoints/Proof/Reductions.lean| below \verb|Words.lean| for the reductions, the row
indices, the bounds of $\mathrm{maxPrimeEven}$, which leave \verb|Params.lean|, and the lemma of
\verb|mod|, which leaves \verb|Tables.lean|; a new file without Mathlib
\verb|Ratpoints/Proof/SetBits.lean|, also below \verb|Words.lean|, for the contract of the trailing
zeros, the walk through the set bits and the class word, which \verb|Words.lean| gives up; the
general facts on words, natural numbers and integers in \verb|Bits.lean|; the rows under a unit and
the rows read in the packing in \verb|Words.lean|; the classes, the offsets, the residues and the
bounds of their products in a new file without Mathlib \verb|Ratpoints/Proof/Packing.lean|, and
above it the bits of an interval with the other range lemmas of the longs in a new file without
Mathlib \verb|Ratpoints/Proof/BitRange.lean|, so that both stay below $1500$ lines; the second stage
and the prime entries in \verb|Stages.lean|; the contracts of the sieve in \verb|Sift.lean|; the
entries of the moduli in \verb|Moduli.lean|, the run in \verb|Denominators.lean|, the shapes in
\verb|Shapes.lean|, the classes' \verb|den_bits| in \verb|TwoAdic.lean|, and the search in
\verb|FindPoints.lean|. Every new module joins its part root and \verb|Ratpoints.lean|. One commit
per item, each passing the project's checks, the constants check and the differential tests, with
the nodes of the earlier chapters that it changes:
\begin{enumerate}
  \item[F3.1] The constants with their facts (Definition~\ref{def:packingConstants}); the high half
    and the reductions through it (Definition~\ref{def:mulHi}, Lemma~\ref{lem:mulHi});
    \verb|mod_mul| and \verb|modLong|'s second range (Definition~\ref{def:modMul},
    Lemma~\ref{lem:modMul}), \verb|mod_mul| not yet called; \verb|Ratpoints/Proof/Reductions.lean|
    without the row indices; the points of a reversed curve, the trailing zeros of a word, the test
    of a candidate, the exact check and the points at infinity to
    \verb|Ratpoints/Model/Candidate.lean| with their guards, about $420$ lines; the nodes of
    Definitions~\ref{def:mulMod}, \ref{def:rbaPack} and~\ref{def:useCLong} and of
    Lemma~\ref{lem:mulMod}, and the files sections of Chapters~\ref{chap:sieve}, \ref{chap:words},
    \ref{chap:setup} and~\ref{chap:stages}.
  \item[F3.2] The fields of the sieve entries and the entry of a modulus with them; the shape of a
    prime entry field by field; the rows under a unit (Lemma~\ref{lem:admits_mul}); the contracts of
    the entry of a modulus and of \verb|sieving_info| with packable entries and tables of inverses,
    with \verb|Ratpoints/Proof/Packing.lean| begun for the packable entries and the bias (importing
    the model of the sieve until F3.3 adds \verb|Ratpoints/Model/Packing.lean|); the nodes of
    Chapters~\ref{chap:tables} and~\ref{chap:stages}.
  \item[F3.3] The classes (Definition~\ref{def:numClasses}, Lemma~\ref{lem:numClasses}),
    \verb|Ratpoints/Model/Packing.lean| and \verb|Ratpoints/Proof/SetBits.lean|, and the classes in
    \verb|Ratpoints/Proof/Packing.lean|; the search takes \verb|den_bits| from the classes; the
    nodes of Definition~\ref{def:twoAdic} and of the search.
  \item[F3.4] The rows of offsets (Definition~\ref{def:classOffsets}, Lemma~\ref{lem:classOffsets}),
    not yet called.
  \item[F3.5] The residues and the bits of an interval (Definitions~\ref{def:fillBpList}
    and~\ref{def:numeratorsToBits}, Lemmas~\ref{lem:fillBpList}, \ref{lem:numeratorsToBits}
    and~\ref{lem:packingLongs}), \verb|Ratpoints/Proof/BitRange.lean|, not yet called.
  \item[F3.6] The packing on: the offsets of the sieve specs, the first two stages with the
    reciprocals, stepped with the sieve specs by two cursors as the C code steps \verb|ssp| and
    \verb|mg|, \verb|small| and the offsets (Lemmas~\ref{lem:reads_row} and~\ref{lem:rowIndices}),
    the step of the walk, the class in the extraction with the base of a bit array stepped before
    each word after the first, as the C code does, the scan and one piece, \verb|set_rows| and the
    bits of one denominator (Definition~\ref{def:siftBits}, Lemma~\ref{lem:siftBits}), one
    denominator with the class, the residues and the offsets, the shapes on the classes, the run
    with the classes and \verb|fill_bp_list|, the search with the rows of offsets
    (Lemma~\ref{lem:soundRun}); comparator; the nodes of Chapters~\ref{chap:sieve},
    \ref{chap:words}, \ref{chap:setup}, \ref{chap:denominators} and~\ref{chap:stages}. One item,
    since every split leaves temporary code that changes the run's contract twice.
  \item[F3.7] The \verb|#guard| examples against the C program's \verb|DEBUG| output; the passages
    of the earlier chapters that F3 makes untrue (Section~\ref{sec:packingEarlier}); the description
    of the repository and \verb|formalization.yaml|.
\end{enumerate}

\section{The evidence}\label{sec:packing-evidence}

\begin{itemize}
  \item A reading of the C code of the packing at the pin, with a second reader who checked it, on
    instrumented builds of the C program (the build of the model with \verb|RATPOINTS_CHUNK=1|, the
    same with composite moduli up to $255$, and the default build with composite moduli up to
    $255$): the classes, \verb|rbainv|, the biases, \verb|dinv|, the offsets, the residues, the bit
    ranges, the masks and the survivors of the first and second stages equal those that the
    definitions give, on $19\,438$ intervals of $600$ cases (above the height $3000$ for the
    denominators up to $40$), each run in the three builds, with $53$ composite moduli, and the
    three branches of \verb|fill_bp_list| on runs with denominators near $2^{24}$, $2^{32}$,
    $10^{12}$ and $2^{62}$; the output of the C program by increasing numerators for each
    denominator, $(a, y)$ before $(a, -y)$, on $1100$ runs; the strides $k \le 3$ on every curve by
    an exhaustive search (all polynomial functions modulo 64 for odd $b$, all truncated reversed
    coefficients for even $b$); the split form of \verb|numerators_to_bits| exact and without
    overflow by a simulation on $64$-bit longs, and no overflow in the extraction. No defect of the
    C code.
  \item A sketch of the model in Lean, which agrees with \verb|ratpoints -q| on the differential
    tests, every point list in its order, and whose classes, offsets, residues, bit ranges, masks
    and test \verb|small| equal those of the instrumented builds on $806$ runs ($51\,584$ classes,
    $25\,233$ sieve entries, $1.14$ million offsets, $7.66$ million residues in all three branches
    with every stride from $2^0$ to $2^3$, $30\,856$ intervals with their pieces and masks, $192$
    pieces outside \verb|small|, $8.39$ million survivors of the first stage and $1.42$ million of
    the second, $25$ composite moduli up to $247$), with a check that two swapped moduli or a
    deleted survivor show as a difference; two reviews of the sketch, each finding checked by a
    second reader, which found no difference from the C code; a timing of the reductions on $64$-bit
    words with the loops stepping two cursors, on one pinned core, by which the model with the
    packing runs \verb|test4high| about a quarter faster than before F3 (1m32s against 2m02s) and
    the default suites about as fast (their time is that of starting the processes).
  \item A prototype of all the proofs on the sketch: the whole library builds, standard axioms,
    every proof at most $20$ lines, the statements that comparator certifies unchanged, the
    differential tests $3315$ compared with none failed and $13$ in \verb|test4high|, and the build
    with other values of the constants ($\mathrm{numStrides}$ $1$ and $8$, $\mathrm{rowBias}$ $0$
    and $2^{30}$, $\mathrm{rbaPack}$ $1$ and $8$ with their shifts); it changed the model of the
    sketch only by the high half of a product, through which the quotient now wraps the reciprocal
    $2^{64}$ of $1$ to $0$, as the C code does, and by the move of the classes, the offsets and the
    residues into a file of their own; the files \verb|Ratpoints/Model/Candidate.lean| and
    \verb|Ratpoints/Proof/BitRange.lean|, the base of the extraction stepped before each word after
    the first, and the names of the row indices and of the range lemmas of the longs, which the
    prototype states under other names, are proposals of this chapter that the prototype does not
    have yet.
\end{itemize}
